% CLASS NOTES BOXES -- standalone example and reusable preamble module.
% Compile this file directly, or load it in another document's preamble:
% \def\NotesBoxesOnly{1}
% % CLASS NOTES BOXES -- standalone example and reusable preamble module.
% Compile this file directly, or load it in another document's preamble:
% \def\NotesBoxesOnly{1}
% % CLASS NOTES BOXES -- standalone example and reusable preamble module.
% Compile this file directly, or load it in another document's preamble:
% \def\NotesBoxesOnly{1}
% % CLASS NOTES BOXES -- standalone example and reusable preamble module.
% Compile this file directly, or load it in another document's preamble:
% \def\NotesBoxesOnly{1}
% \input{class_notes_boxes.tex}
% See class_notes_boxes_guide.md for the complete catalog and instructions.

\ifdefined\NotesBoxesOnly
\else
  \documentclass[11pt]{article}
  \usepackage[margin=1in]{geometry}
  \usepackage[T1]{fontenc}
  \usepackage{lmodern}
\fi

\usepackage{amsmath,amssymb}
\usepackage[most]{tcolorbox}
\usepackage{graphicx}
\usepackage{booktabs}
\usepackage{listings}

% Shared appearance. The title text is always bold; body text is normal
% unless an environment (answer) explicitly changes it.
\tcbset{
  notesbase/.style={enhanced,breakable,boxrule=0.7pt,arc=2mm,
    left=3mm,right=3mm,top=2mm,bottom=2mm,
    before skip=10pt,after skip=10pt,
    fonttitle=\bfseries,fontupper=\normalfont,
    coltext=black,coltitle=white},
  notesblue/.style={colback=blue!3!white,colframe=blue!45!black,
    colbacktitle=blue!45!black},
  notesgreen/.style={colback=green!6!white,colframe=green!30!black,
    colbacktitle=green!12!white,coltitle=green!20!black},
  notesgray/.style={colback=black!2!white,colframe=black!55,
    colbacktitle=black!65},
  notesamber/.style={colback=yellow!12!white,colframe=orange!60!black,
    colbacktitle=yellow!28!white,coltitle=black!85},
  notespurple/.style={colback=violet!4!white,colframe=violet!45!black,
    colbacktitle=violet!45!black},
  noteswarning/.style={colback=black!85,colframe=yellow!80!orange,
    colbacktitle=black!90,coltitle=yellow!85!white,coltext=yellow!85!white},
  notesred/.style={colback=red!5!white,colframe=red!65!black,
    colbacktitle=red!65!black}
}

% Convenient optional fields for any box.
\newcommand{\NotesField}[2]{\par\noindent\textbf{#1:} #2\par}

% A literal code style: spaces, braces, underscores and percent signs
% inside lstlisting do not need LaTeX escaping.
\lstdefinestyle{notescode}{basicstyle=\small\ttfamily,
  breaklines=true,columns=fullflexible,keepspaces=true,
  showstringspaces=false,tabsize=4}

% BOX DEFINITIONS. All boxes accept optional tcolorbox keys in [...].

% Everyday learning
\newtcolorbox{note}[1][]{notesbase,notesgreen,
  title={Note},#1}
\newtcolorbox{definitionbox}[1][]{notesbase,notesblue,
  title={Definition},#1}
\newtcolorbox{conceptbox}[1][]{notesbase,notesblue,
  title={Key Concept},#1}
\newtcolorbox{notationbox}[1][]{notesbase,notesblue,
  title={Notation},#1}
\newtcolorbox{objectivebox}[1][]{notesbase,notesgreen,
  title={Learning Objectives},#1}
\newtcolorbox{contextbox}[1][]{notesbase,notesgray,
  title={Background / Context},#1}
\newtcolorbox{connectionbox}[1][]{notesbase,notesgreen,
  title={Connection},#1}
\newtcolorbox{summarybox}[1][]{notesbase,notesgreen,
  title={Summary},#1}

% Questions and study
\newtcolorbox[auto counter]{question}[1][]{notesbase,notesblue,
  title={Question~\thetcbcounter},#1}
\newtcolorbox{answer}[1][]{notesbase,notesgray,
  title={Answer},fontupper=\itshape,#1}
\newtcolorbox{examplebox}[1][]{notesbase,notesgray,
  title={Example},#1}
\newtcolorbox{solutionbox}[1][]{notesbase,notesgray,
  title={Worked Solution},#1}
\newtcolorbox{practicebox}[1][]{notesbase,notesblue,
  title={Practice},#1}
\newtcolorbox{recallbox}[1][]{notesbase,notespurple,
  title={Active Recall},#1}
\newtcolorbox{hintbox}[1][]{notesbase,notesgreen,
  title={Hint},#1}
\newtcolorbox{misconceptionbox}[1][]{notesbase,notesamber,
  title={Common Mistake},#1}
\newtcolorbox{askbox}[1][]{notesbase,notesamber,
  title={Ask / Clarify},#1}
\newtcolorbox{actionbox}[1][]{notesbase,notespurple,
  title={Action Item},#1}
\newtcolorbox{reflectionbox}[1][]{notesbase,notespurple,
  title={Reflection},#1}

% Math and quantitative reasoning
\newtcolorbox{formulabox}[1][]{notesbase,notesblue,
  title={Formula / Rule},#1}
\newtcolorbox{theorembox}[1][]{notesbase,notesblue,
  title={Theorem / Principle},#1}
\newtcolorbox{proofbox}[1][]{notesbase,notesgray,
  title={Proof},#1}
\newtcolorbox{derivationbox}[1][]{notesbase,notesgray,
  title={Derivation},#1}
\newtcolorbox{assumptionbox}[1][]{notesbase,notesamber,
  title={Assumptions},#1}
\newtcolorbox{givenbox}[1][]{notesbase,notesgray,
  title={Given / Find},#1}
\newtcolorbox{unitsbox}[1][]{notesbase,notesblue,
  title={Units / Dimensions},#1}
\newtcolorbox{checkbox}[1][]{notesbase,notesgreen,
  title={Result Check},#1}
\newtcolorbox{specialcasebox}[1][]{notesbase,notesamber,
  title={Special Case / Exception},#1}

% Science, labs, and research
\newtcolorbox{hypothesisbox}[1][]{notesbase,notesblue,
  title={Hypothesis},#1}
\newtcolorbox{variablesbox}[1][]{notesbase,notesblue,
  title={Variables / Controls},#1}
\newtcolorbox{equipmentbox}[1][]{notesbase,notesgray,
  title={Equipment / Materials},#1}
\newtcolorbox{procedurebox}[1][]{notesbase,notesgray,
  title={Procedure / Method},#1}
\newtcolorbox{observationbox}[1][]{notesbase,notesgray,
  title={Observation},#1}
\newtcolorbox{databox}[1][]{notesbase,notesgray,
  title={Data},#1}
\newtcolorbox{uncertaintybox}[1][]{notesbase,notesamber,
  title={Uncertainty / Error},#1}
\newtcolorbox{analysisbox}[1][]{notesbase,notesgray,
  title={Analysis},#1}
\newtcolorbox{conclusionbox}[1][]{notesbase,notesgreen,
  title={Conclusion},#1}

% Humanities, languages, and social sciences
\newtcolorbox{claimbox}[1][]{notesbase,notesblue,
  title={Claim / Thesis},#1}
\newtcolorbox{evidencebox}[1][]{notesbase,notesgray,
  title={Evidence / Source},#1}
\newtcolorbox{quotationbox}[1][]{notesbase,notesgray,
  title={Quotation},#1}
\newtcolorbox{interpretationbox}[1][]{notesbase,notespurple,
  title={Interpretation},#1}
\newtcolorbox{counterargumentbox}[1][]{notesbase,notesamber,
  title={Counterargument},#1}
\newtcolorbox{comparisonbox}[1][]{notesbase,notesgray,
  title={Comparison},#1}
\newtcolorbox{timelinebox}[1][]{notesbase,notesgray,
  title={Timeline},#1}
\newtcolorbox{casebox}[1][]{notesbase,notesgray,
  title={Case Study},#1}
\newtcolorbox{vocabularybox}[1][]{notesbase,notesblue,
  title={Vocabulary},#1}
\newtcolorbox{grammarbox}[1][]{notesbase,notesblue,
  title={Grammar / Language Pattern},#1}

% Programming, engineering, and practical work
\newtcolorbox{codebox}[1][]{notesbase,notesgray,
  title={Code},#1}
\newtcolorbox{outputbox}[1][]{notesbase,notesgray,
  title={Expected / Actual Output},#1}
\newtcolorbox{debugbox}[1][]{notesbase,notesamber,
  title={Error / Fix},#1}
\newtcolorbox{designbox}[1][]{notesbase,notespurple,
  title={Design Decision},#1}
\newtcolorbox{testbox}[1][]{notesbase,notesgreen,
  title={Test Case},#1}
\newtcolorbox{figurebox}[1][]{notesbase,notesgray,
  title={Figure / Diagram},#1}
\newtcolorbox{techniquebox}[1][]{notesbase,notesgray,
  title={Technique / Skill},#1}
\newtcolorbox{feedbackbox}[1][]{notesbase,notespurple,
  title={Feedback / Critique},#1}

% Safety and cautions
\newtcolorbox{warning}[1][]{notesbase,noteswarning,
  title={Warning},#1}
\newtcolorbox{danger}[1][]{notesbase,notesred,
  title={Danger},#1}

% DEMONSTRATION DOCUMENT: skipped when NotesBoxesOnly is defined.
\ifdefined\NotesBoxesOnly
\else
\usepackage[hidelinks]{hyperref}
\title{Class Notes Box Collection}
\author{Reusable LaTeX environments for learning and review}
\date{}
\begin{document}
\maketitle
This file contains 57 reusable box types for a broad range of classes.
The examples demonstrate formatting; replace their contents with your notes.
Use ordinary paragraphs between boxes to keep your notes readable.
\tableofcontents
\clearpage
\section{Everyday learning}
\subsection{Note}
{\small\texttt{note}: Supplementary explanation or a helpful reminder.}
\begin{note}
Review this topic after the next lecture.
\end{note}
\subsection{Definition}
{\small\texttt{definitionbox}: A term with its precise meaning.}
\begin{definitionbox}
Velocity is the rate of change of position with respect to time.
\end{definitionbox}
\subsection{Key Concept}
{\small\texttt{conceptbox}: An important idea in your own words.}
\begin{conceptbox}
A net force changes the velocity of an object.
\end{conceptbox}
\subsection{Notation}
{\small\texttt{notationbox}: Explain symbols, abbreviations, and conventions.}
\begin{notationbox}
$m$ denotes mass; $a$ denotes acceleration.
\end{notationbox}
\subsection{Learning Objectives}
{\small\texttt{objectivebox}: State what you should be able to explain or do.}
\begin{objectivebox}
Explain the difference between speed and velocity.
\end{objectivebox}
\subsection{Background / Context}
{\small\texttt{contextbox}: Give background needed to understand a topic.}
\begin{contextbox}
This lesson assumes familiarity with vectors.
\end{contextbox}
\subsection{Connection}
{\small\texttt{connectionbox}: Connect a topic to earlier lessons or another subject.}
\begin{connectionbox}
Velocity connects calculus derivatives with physical motion.
\end{connectionbox}
\subsection{Summary}
{\small\texttt{summarybox}: Record the main takeaways of a lesson.}
\begin{summarybox}
Identify forces, choose axes, and then apply the equations of motion.
\end{summarybox}

\clearpage
\section{Questions and study}
\subsection{Question}
{\small\texttt{question}: A numbered question; numbering increases automatically.}
\begin{question}
A net force of $20\,\mathrm{N}$ acts on a $5\,\mathrm{kg}$ object. Find its acceleration.
\end{question}
\subsection{Answer}
{\small\texttt{answer}: An answer with automatically italicized prose.}
\begin{answer}
$a=F_{\mathrm{net}}/m=4\,\mathrm{m/s^2}$.
\end{answer}
\subsection{Example}
{\small\texttt{examplebox}: A concrete illustration of an idea.}
\begin{examplebox}
A car travelling at constant speed around a bend still accelerates.
\end{examplebox}
\subsection{Worked Solution}
{\small\texttt{solutionbox}: Show the reasoning and steps of a calculation.}
\begin{solutionbox}
Choose the positive direction. Calculate the net force. Divide by mass.
\end{solutionbox}
\subsection{Practice}
{\small\texttt{practicebox}: An exercise to complete independently.}
\begin{practicebox}
Repeat the calculation with twice the mass.
\end{practicebox}
\subsection{Active Recall}
{\small\texttt{recallbox}: A prompt to answer without consulting the notes.}
\begin{recallbox}
State the three laws of motion in your own words.
\end{recallbox}
\subsection{Hint}
{\small\texttt{hintbox}: A small clue that does not reveal the entire solution.}
\begin{hintbox}
Start by drawing a free-body diagram.
\end{hintbox}
\subsection{Common Mistake}
{\small\texttt{misconceptionbox}: Identify an error, explain why it fails, and give a correction.}
\begin{misconceptionbox}
Constant speed does not always mean zero acceleration.
\end{misconceptionbox}
\subsection{Ask / Clarify}
{\small\texttt{askbox}: An unanswered question for an instructor or later research.}
\begin{askbox}
Under what conditions can friction be ignored?
\end{askbox}
\subsection{Action Item}
{\small\texttt{actionbox}: A task, its next step, and an optional due date.}
\begin{actionbox}
Complete practice problems 1--5 before the next class.
\end{actionbox}
\subsection{Reflection}
{\small\texttt{reflectionbox}: Evaluate your understanding, approach, or learning progress.}
\begin{reflectionbox}
I can draw the forces, but I need more practice choosing components.
\end{reflectionbox}

\clearpage
\section{Math and quantitative reasoning}
\subsection{Formula / Rule}
{\small\texttt{formulabox}: An equation or rule, with variables, units, and conditions of use.}
\begin{formulabox}
$F_{\mathrm{net}}=ma$ for constant mass in an inertial reference frame.
\end{formulabox}
\subsection{Theorem / Principle}
{\small\texttt{theorembox}: A formal result together with its hypotheses.}
\begin{theorembox}
If a function is continuous on a closed interval, it attains a maximum and a minimum.
\end{theorembox}
\subsection{Proof}
{\small\texttt{proofbox}: A logical argument establishing a claim.}
\begin{proofbox}
State the assumptions, justify each inference, and establish the conclusion.
\end{proofbox}
\subsection{Derivation}
{\small\texttt{derivationbox}: Develop an equation from definitions or established relationships.}
\begin{derivationbox}
Starting with $v=dx/dt$, integrate to relate displacement to velocity.
\end{derivationbox}
\subsection{Assumptions}
{\small\texttt{assumptionbox}: Make the conditions of a model or calculation explicit.}
\begin{assumptionbox}
Neglect air resistance and treat gravitational acceleration as constant.
\end{assumptionbox}
\subsection{Given / Find}
{\small\texttt{givenbox}: Separate known quantities from the target quantity.}
\begin{givenbox}
Given: mass and net force. Find: acceleration.
\end{givenbox}
\subsection{Units / Dimensions}
{\small\texttt{unitsbox}: Define units, convert them, or check dimensions.}
\begin{unitsbox}
$1\,\mathrm{N}=1\,\mathrm{kg\,m/s^2}$.
\end{unitsbox}
\subsection{Result Check}
{\small\texttt{checkbox}: Check sign, magnitude, units, limiting behavior, or plausibility.}
\begin{checkbox}
Doubling the mass at fixed force should halve the acceleration.
\end{checkbox}
\subsection{Special Case / Exception}
{\small\texttt{specialcasebox}: Identify boundaries, exceptions, or cases requiring different treatment.}
\begin{specialcasebox}
The expression $1/x$ is undefined at $x=0$.
\end{specialcasebox}

\clearpage
\section{Science, labs, and research}
\subsection{Hypothesis}
{\small\texttt{hypothesisbox}: A testable prediction with a rationale.}
\begin{hypothesisbox}
Increasing pendulum length will increase its period for small oscillations.
\end{hypothesisbox}
\subsection{Variables / Controls}
{\small\texttt{variablesbox}: List independent, dependent, and controlled variables.}
\begin{variablesbox}
Independent: length. Dependent: period. Controlled: release angle.
\end{variablesbox}
\subsection{Equipment / Materials}
{\small\texttt{equipmentbox}: List apparatus, quantities, and relevant specifications.}
\begin{equipmentbox}
String, a pendulum bob, a meter ruler, and a stopwatch.
\end{equipmentbox}
\subsection{Procedure / Method}
{\small\texttt{procedurebox}: Numbered steps for carrying out a repeatable task.}
\begin{procedurebox}
Measure the length, time several oscillations, and repeat the measurement.
\end{procedurebox}
\subsection{Observation}
{\small\texttt{observationbox}: Describe what was observed without mixing in interpretation.}
\begin{observationbox}
The measured period increased when the string was lengthened.
\end{observationbox}
\subsection{Data}
{\small\texttt{databox}: Measurements or a small table with units and labels.}
\begin{databox}
Record the measured length and period for every trial.
\end{databox}
\subsection{Uncertainty / Error}
{\small\texttt{uncertaintybox}: Discuss measurement uncertainty and systematic or random errors.}
\begin{uncertaintybox}
Reaction time contributes uncertainty to stopwatch measurements.
\end{uncertaintybox}
\subsection{Analysis}
{\small\texttt{analysisbox}: Explain how data or evidence supports your interpretation.}
\begin{analysisbox}
Compare the trend with the model prediction before drawing a conclusion.
\end{analysisbox}
\subsection{Conclusion}
{\small\texttt{conclusionbox}: State the result, evidence, and limits of the conclusion.}
\begin{conclusionbox}
The observed trend supports the prediction within the tested range.
\end{conclusionbox}

\clearpage
\section{Humanities, languages, and social sciences}
\subsection{Claim / Thesis}
{\small\texttt{claimbox}: A position that needs support from evidence.}
\begin{claimbox}
The author uses setting to establish the central conflict.
\end{claimbox}
\subsection{Evidence / Source}
{\small\texttt{evidencebox}: Evidence with enough source information to find it again.}
\begin{evidencebox}
Record the author, work, page or section, and the supporting detail.
\end{evidencebox}
\subsection{Quotation}
{\small\texttt{quotationbox}: An exact quotation with its source and page or location.}
\begin{quotationbox}
Place a short exact quotation here and record its source.
\end{quotationbox}
\subsection{Interpretation}
{\small\texttt{interpretationbox}: Explain what a passage, event, or piece of evidence means.}
\begin{interpretationbox}
Explain how the quoted language supports the claim.
\end{interpretationbox}
\subsection{Counterargument}
{\small\texttt{counterargumentbox}: Present an objection fairly and evaluate a response.}
\begin{counterargumentbox}
An alternative interpretation may explain the same evidence differently.
\end{counterargumentbox}
\subsection{Comparison}
{\small\texttt{comparisonbox}: Compare related ideas using the same criteria.}
\begin{comparisonbox}
Compare two theories by assumptions, predictions, evidence, and limitations.
\end{comparisonbox}
\subsection{Timeline}
{\small\texttt{timelinebox}: Events in chronological order with their significance.}
\begin{timelinebox}
Record each date, event, and why it matters.
\end{timelinebox}
\subsection{Case Study}
{\small\texttt{casebox}: Apply a model to a specific person, organization, event, or scenario.}
\begin{casebox}
Describe the situation, identify the relevant theory, and assess its fit.
\end{casebox}
\subsection{Vocabulary}
{\small\texttt{vocabularybox}: A word, meaning, pronunciation, and an example sentence.}
\begin{vocabularybox}
Word: concise. Meaning: brief but clear. Example: Write a concise summary.
\end{vocabularybox}
\subsection{Grammar / Language Pattern}
{\small\texttt{grammarbox}: A language pattern with examples and exceptions.}
\begin{grammarbox}
Record the rule, an example, and an exception if one applies.
\end{grammarbox}

\clearpage
\section{Programming, engineering, and practical work}
\subsection{Code}
{\small\texttt{codebox}: A short code listing with its language and purpose.}
\begin{codebox}
Use a lstlisting environment inside this box for literal code.
\end{codebox}
\subsection{Expected / Actual Output}
{\small\texttt{outputbox}: Record a program result, calculation output, or observed behavior.}
\begin{outputbox}
Expected output: 4.0. Actual output: 4.0.
\end{outputbox}
\subsection{Error / Fix}
{\small\texttt{debugbox}: Record the symptom, root cause, correction, and verification.}
\begin{debugbox}
A null value caused the error; validate input before dereferencing it.
\end{debugbox}
\subsection{Design Decision}
{\small\texttt{designbox}: Record a decision, alternatives, reasons, and trade-offs.}
\begin{designbox}
Separate parsing from rendering so either can change independently.
\end{designbox}
\subsection{Test Case}
{\small\texttt{testbox}: Specify input, expected result, and pass or fail evidence.}
\begin{testbox}
Input: zero net force. Expected acceleration: zero.
\end{testbox}
\subsection{Figure / Diagram}
{\small\texttt{figurebox}: A labeled diagram or image with a descriptive caption.}
\begin{figurebox}
Insert a diagram and explain what its arrows, axes, or labels represent.
\end{figurebox}
\subsection{Technique / Skill}
{\small\texttt{techniquebox}: Describe a practical technique and its success criteria.}
\begin{techniquebox}
Record the setup, movement or action, and signs of correct execution.
\end{techniquebox}
\subsection{Feedback / Critique}
{\small\texttt{feedbackbox}: Record feedback and a concrete improvement to make.}
\begin{feedbackbox}
Feedback: explain assumptions. Next revision: add an assumptions box.
\end{feedbackbox}

\clearpage
\section{Safety and cautions}
\subsection{Warning}
{\small\texttt{warning}: A caution requiring attention; yellow text on a dark background.}
\begin{warning}
Verify the operating limits before using the equipment.
\end{warning}
\subsection{Danger}
{\small\texttt{danger}: A serious hazard; red frame and title bar with a pale red body.}
\begin{danger}
Disconnect power before touching exposed electrical components.
\end{danger}

\clearpage
\section{Richer contents and customization}
\begin{formulabox}[title={Newton's Second Law},label={box:newton}]
\[
  \sum \vec{F}=m\vec{a}
\]
\NotesField{Symbols}{$m$ is mass; $\vec{a}$ is acceleration.}
\NotesField{Units}{Force in newtons, mass in kilograms, acceleration in $\mathrm{m/s^2}$.}
\NotesField{When it applies}{Constant mass, in an inertial reference frame.}
\end{formulabox}

\begin{procedurebox}
\begin{enumerate}
  \item Draw and label the forces.
  \item Choose coordinate axes.
  \item Resolve forces into components.
  \item Apply the equation along each axis.
\end{enumerate}
\end{procedurebox}

\begin{databox}[title={Measurements}]
\centering
\begin{tabular}{rr}
\toprule
Mass ($\mathrm{kg}$) & Force ($\mathrm{N}$) \\
\midrule
2 & 8 \\
5 & 20 \\
\bottomrule
\end{tabular}
\end{databox}

\begin{codebox}[title={Python: acceleration}]
\begin{lstlisting}[style=notescode,language=Python]
def acceleration(net_force, mass):
    if mass <= 0:
        raise ValueError("mass must be positive")
    return net_force / mass

print(acceleration(20, 5))
\end{lstlisting}
\end{codebox}

\begin{outputbox}
\begin{lstlisting}[style=notescode]
4.0
\end{lstlisting}
\end{outputbox}

\begin{question}[title={Question 12: custom display title}]
This changes only the displayed title, not the automatic counter.
\end{question}
\begin{answer}
The prose is italic; mathematical notation such as $F=ma$ keeps its normal mathematical formatting.
\end{answer}

\begin{note}[title={A customized title},colback=green!10!white]
Use optional keys to change one box without changing all boxes.
\end{note}

\end{document}
\fi

% See class_notes_boxes_guide.md for the complete catalog and instructions.

\ifdefined\NotesBoxesOnly
\else
  \documentclass[11pt]{article}
  \usepackage[margin=1in]{geometry}
  \usepackage[T1]{fontenc}
  \usepackage{lmodern}
\fi

\usepackage{amsmath,amssymb}
\usepackage[most]{tcolorbox}
\usepackage{graphicx}
\usepackage{booktabs}
\usepackage{listings}

% Shared appearance. The title text is always bold; body text is normal
% unless an environment (answer) explicitly changes it.
\tcbset{
  notesbase/.style={enhanced,breakable,boxrule=0.7pt,arc=2mm,
    left=3mm,right=3mm,top=2mm,bottom=2mm,
    before skip=10pt,after skip=10pt,
    fonttitle=\bfseries,fontupper=\normalfont,
    coltext=black,coltitle=white},
  notesblue/.style={colback=blue!3!white,colframe=blue!45!black,
    colbacktitle=blue!45!black},
  notesgreen/.style={colback=green!6!white,colframe=green!30!black,
    colbacktitle=green!12!white,coltitle=green!20!black},
  notesgray/.style={colback=black!2!white,colframe=black!55,
    colbacktitle=black!65},
  notesamber/.style={colback=yellow!12!white,colframe=orange!60!black,
    colbacktitle=yellow!28!white,coltitle=black!85},
  notespurple/.style={colback=violet!4!white,colframe=violet!45!black,
    colbacktitle=violet!45!black},
  noteswarning/.style={colback=black!85,colframe=yellow!80!orange,
    colbacktitle=black!90,coltitle=yellow!85!white,coltext=yellow!85!white},
  notesred/.style={colback=red!5!white,colframe=red!65!black,
    colbacktitle=red!65!black}
}

% Convenient optional fields for any box.
\newcommand{\NotesField}[2]{\par\noindent\textbf{#1:} #2\par}

% A literal code style: spaces, braces, underscores and percent signs
% inside lstlisting do not need LaTeX escaping.
\lstdefinestyle{notescode}{basicstyle=\small\ttfamily,
  breaklines=true,columns=fullflexible,keepspaces=true,
  showstringspaces=false,tabsize=4}

% BOX DEFINITIONS. All boxes accept optional tcolorbox keys in [...].

% Everyday learning
\newtcolorbox{note}[1][]{notesbase,notesgreen,
  title={Note},#1}
\newtcolorbox{definitionbox}[1][]{notesbase,notesblue,
  title={Definition},#1}
\newtcolorbox{conceptbox}[1][]{notesbase,notesblue,
  title={Key Concept},#1}
\newtcolorbox{notationbox}[1][]{notesbase,notesblue,
  title={Notation},#1}
\newtcolorbox{objectivebox}[1][]{notesbase,notesgreen,
  title={Learning Objectives},#1}
\newtcolorbox{contextbox}[1][]{notesbase,notesgray,
  title={Background / Context},#1}
\newtcolorbox{connectionbox}[1][]{notesbase,notesgreen,
  title={Connection},#1}
\newtcolorbox{summarybox}[1][]{notesbase,notesgreen,
  title={Summary},#1}

% Questions and study
\newtcolorbox[auto counter]{question}[1][]{notesbase,notesblue,
  title={Question~\thetcbcounter},#1}
\newtcolorbox{answer}[1][]{notesbase,notesgray,
  title={Answer},fontupper=\itshape,#1}
\newtcolorbox{examplebox}[1][]{notesbase,notesgray,
  title={Example},#1}
\newtcolorbox{solutionbox}[1][]{notesbase,notesgray,
  title={Worked Solution},#1}
\newtcolorbox{practicebox}[1][]{notesbase,notesblue,
  title={Practice},#1}
\newtcolorbox{recallbox}[1][]{notesbase,notespurple,
  title={Active Recall},#1}
\newtcolorbox{hintbox}[1][]{notesbase,notesgreen,
  title={Hint},#1}
\newtcolorbox{misconceptionbox}[1][]{notesbase,notesamber,
  title={Common Mistake},#1}
\newtcolorbox{askbox}[1][]{notesbase,notesamber,
  title={Ask / Clarify},#1}
\newtcolorbox{actionbox}[1][]{notesbase,notespurple,
  title={Action Item},#1}
\newtcolorbox{reflectionbox}[1][]{notesbase,notespurple,
  title={Reflection},#1}

% Math and quantitative reasoning
\newtcolorbox{formulabox}[1][]{notesbase,notesblue,
  title={Formula / Rule},#1}
\newtcolorbox{theorembox}[1][]{notesbase,notesblue,
  title={Theorem / Principle},#1}
\newtcolorbox{proofbox}[1][]{notesbase,notesgray,
  title={Proof},#1}
\newtcolorbox{derivationbox}[1][]{notesbase,notesgray,
  title={Derivation},#1}
\newtcolorbox{assumptionbox}[1][]{notesbase,notesamber,
  title={Assumptions},#1}
\newtcolorbox{givenbox}[1][]{notesbase,notesgray,
  title={Given / Find},#1}
\newtcolorbox{unitsbox}[1][]{notesbase,notesblue,
  title={Units / Dimensions},#1}
\newtcolorbox{checkbox}[1][]{notesbase,notesgreen,
  title={Result Check},#1}
\newtcolorbox{specialcasebox}[1][]{notesbase,notesamber,
  title={Special Case / Exception},#1}

% Science, labs, and research
\newtcolorbox{hypothesisbox}[1][]{notesbase,notesblue,
  title={Hypothesis},#1}
\newtcolorbox{variablesbox}[1][]{notesbase,notesblue,
  title={Variables / Controls},#1}
\newtcolorbox{equipmentbox}[1][]{notesbase,notesgray,
  title={Equipment / Materials},#1}
\newtcolorbox{procedurebox}[1][]{notesbase,notesgray,
  title={Procedure / Method},#1}
\newtcolorbox{observationbox}[1][]{notesbase,notesgray,
  title={Observation},#1}
\newtcolorbox{databox}[1][]{notesbase,notesgray,
  title={Data},#1}
\newtcolorbox{uncertaintybox}[1][]{notesbase,notesamber,
  title={Uncertainty / Error},#1}
\newtcolorbox{analysisbox}[1][]{notesbase,notesgray,
  title={Analysis},#1}
\newtcolorbox{conclusionbox}[1][]{notesbase,notesgreen,
  title={Conclusion},#1}

% Humanities, languages, and social sciences
\newtcolorbox{claimbox}[1][]{notesbase,notesblue,
  title={Claim / Thesis},#1}
\newtcolorbox{evidencebox}[1][]{notesbase,notesgray,
  title={Evidence / Source},#1}
\newtcolorbox{quotationbox}[1][]{notesbase,notesgray,
  title={Quotation},#1}
\newtcolorbox{interpretationbox}[1][]{notesbase,notespurple,
  title={Interpretation},#1}
\newtcolorbox{counterargumentbox}[1][]{notesbase,notesamber,
  title={Counterargument},#1}
\newtcolorbox{comparisonbox}[1][]{notesbase,notesgray,
  title={Comparison},#1}
\newtcolorbox{timelinebox}[1][]{notesbase,notesgray,
  title={Timeline},#1}
\newtcolorbox{casebox}[1][]{notesbase,notesgray,
  title={Case Study},#1}
\newtcolorbox{vocabularybox}[1][]{notesbase,notesblue,
  title={Vocabulary},#1}
\newtcolorbox{grammarbox}[1][]{notesbase,notesblue,
  title={Grammar / Language Pattern},#1}

% Programming, engineering, and practical work
\newtcolorbox{codebox}[1][]{notesbase,notesgray,
  title={Code},#1}
\newtcolorbox{outputbox}[1][]{notesbase,notesgray,
  title={Expected / Actual Output},#1}
\newtcolorbox{debugbox}[1][]{notesbase,notesamber,
  title={Error / Fix},#1}
\newtcolorbox{designbox}[1][]{notesbase,notespurple,
  title={Design Decision},#1}
\newtcolorbox{testbox}[1][]{notesbase,notesgreen,
  title={Test Case},#1}
\newtcolorbox{figurebox}[1][]{notesbase,notesgray,
  title={Figure / Diagram},#1}
\newtcolorbox{techniquebox}[1][]{notesbase,notesgray,
  title={Technique / Skill},#1}
\newtcolorbox{feedbackbox}[1][]{notesbase,notespurple,
  title={Feedback / Critique},#1}

% Safety and cautions
\newtcolorbox{warning}[1][]{notesbase,noteswarning,
  title={Warning},#1}
\newtcolorbox{danger}[1][]{notesbase,notesred,
  title={Danger},#1}

% DEMONSTRATION DOCUMENT: skipped when NotesBoxesOnly is defined.
\ifdefined\NotesBoxesOnly
\else
\usepackage[hidelinks]{hyperref}
\title{Class Notes Box Collection}
\author{Reusable LaTeX environments for learning and review}
\date{}
\begin{document}
\maketitle
This file contains 57 reusable box types for a broad range of classes.
The examples demonstrate formatting; replace their contents with your notes.
Use ordinary paragraphs between boxes to keep your notes readable.
\tableofcontents
\clearpage
\section{Everyday learning}
\subsection{Note}
{\small\texttt{note}: Supplementary explanation or a helpful reminder.}
\begin{note}
Review this topic after the next lecture.
\end{note}
\subsection{Definition}
{\small\texttt{definitionbox}: A term with its precise meaning.}
\begin{definitionbox}
Velocity is the rate of change of position with respect to time.
\end{definitionbox}
\subsection{Key Concept}
{\small\texttt{conceptbox}: An important idea in your own words.}
\begin{conceptbox}
A net force changes the velocity of an object.
\end{conceptbox}
\subsection{Notation}
{\small\texttt{notationbox}: Explain symbols, abbreviations, and conventions.}
\begin{notationbox}
$m$ denotes mass; $a$ denotes acceleration.
\end{notationbox}
\subsection{Learning Objectives}
{\small\texttt{objectivebox}: State what you should be able to explain or do.}
\begin{objectivebox}
Explain the difference between speed and velocity.
\end{objectivebox}
\subsection{Background / Context}
{\small\texttt{contextbox}: Give background needed to understand a topic.}
\begin{contextbox}
This lesson assumes familiarity with vectors.
\end{contextbox}
\subsection{Connection}
{\small\texttt{connectionbox}: Connect a topic to earlier lessons or another subject.}
\begin{connectionbox}
Velocity connects calculus derivatives with physical motion.
\end{connectionbox}
\subsection{Summary}
{\small\texttt{summarybox}: Record the main takeaways of a lesson.}
\begin{summarybox}
Identify forces, choose axes, and then apply the equations of motion.
\end{summarybox}

\clearpage
\section{Questions and study}
\subsection{Question}
{\small\texttt{question}: A numbered question; numbering increases automatically.}
\begin{question}
A net force of $20\,\mathrm{N}$ acts on a $5\,\mathrm{kg}$ object. Find its acceleration.
\end{question}
\subsection{Answer}
{\small\texttt{answer}: An answer with automatically italicized prose.}
\begin{answer}
$a=F_{\mathrm{net}}/m=4\,\mathrm{m/s^2}$.
\end{answer}
\subsection{Example}
{\small\texttt{examplebox}: A concrete illustration of an idea.}
\begin{examplebox}
A car travelling at constant speed around a bend still accelerates.
\end{examplebox}
\subsection{Worked Solution}
{\small\texttt{solutionbox}: Show the reasoning and steps of a calculation.}
\begin{solutionbox}
Choose the positive direction. Calculate the net force. Divide by mass.
\end{solutionbox}
\subsection{Practice}
{\small\texttt{practicebox}: An exercise to complete independently.}
\begin{practicebox}
Repeat the calculation with twice the mass.
\end{practicebox}
\subsection{Active Recall}
{\small\texttt{recallbox}: A prompt to answer without consulting the notes.}
\begin{recallbox}
State the three laws of motion in your own words.
\end{recallbox}
\subsection{Hint}
{\small\texttt{hintbox}: A small clue that does not reveal the entire solution.}
\begin{hintbox}
Start by drawing a free-body diagram.
\end{hintbox}
\subsection{Common Mistake}
{\small\texttt{misconceptionbox}: Identify an error, explain why it fails, and give a correction.}
\begin{misconceptionbox}
Constant speed does not always mean zero acceleration.
\end{misconceptionbox}
\subsection{Ask / Clarify}
{\small\texttt{askbox}: An unanswered question for an instructor or later research.}
\begin{askbox}
Under what conditions can friction be ignored?
\end{askbox}
\subsection{Action Item}
{\small\texttt{actionbox}: A task, its next step, and an optional due date.}
\begin{actionbox}
Complete practice problems 1--5 before the next class.
\end{actionbox}
\subsection{Reflection}
{\small\texttt{reflectionbox}: Evaluate your understanding, approach, or learning progress.}
\begin{reflectionbox}
I can draw the forces, but I need more practice choosing components.
\end{reflectionbox}

\clearpage
\section{Math and quantitative reasoning}
\subsection{Formula / Rule}
{\small\texttt{formulabox}: An equation or rule, with variables, units, and conditions of use.}
\begin{formulabox}
$F_{\mathrm{net}}=ma$ for constant mass in an inertial reference frame.
\end{formulabox}
\subsection{Theorem / Principle}
{\small\texttt{theorembox}: A formal result together with its hypotheses.}
\begin{theorembox}
If a function is continuous on a closed interval, it attains a maximum and a minimum.
\end{theorembox}
\subsection{Proof}
{\small\texttt{proofbox}: A logical argument establishing a claim.}
\begin{proofbox}
State the assumptions, justify each inference, and establish the conclusion.
\end{proofbox}
\subsection{Derivation}
{\small\texttt{derivationbox}: Develop an equation from definitions or established relationships.}
\begin{derivationbox}
Starting with $v=dx/dt$, integrate to relate displacement to velocity.
\end{derivationbox}
\subsection{Assumptions}
{\small\texttt{assumptionbox}: Make the conditions of a model or calculation explicit.}
\begin{assumptionbox}
Neglect air resistance and treat gravitational acceleration as constant.
\end{assumptionbox}
\subsection{Given / Find}
{\small\texttt{givenbox}: Separate known quantities from the target quantity.}
\begin{givenbox}
Given: mass and net force. Find: acceleration.
\end{givenbox}
\subsection{Units / Dimensions}
{\small\texttt{unitsbox}: Define units, convert them, or check dimensions.}
\begin{unitsbox}
$1\,\mathrm{N}=1\,\mathrm{kg\,m/s^2}$.
\end{unitsbox}
\subsection{Result Check}
{\small\texttt{checkbox}: Check sign, magnitude, units, limiting behavior, or plausibility.}
\begin{checkbox}
Doubling the mass at fixed force should halve the acceleration.
\end{checkbox}
\subsection{Special Case / Exception}
{\small\texttt{specialcasebox}: Identify boundaries, exceptions, or cases requiring different treatment.}
\begin{specialcasebox}
The expression $1/x$ is undefined at $x=0$.
\end{specialcasebox}

\clearpage
\section{Science, labs, and research}
\subsection{Hypothesis}
{\small\texttt{hypothesisbox}: A testable prediction with a rationale.}
\begin{hypothesisbox}
Increasing pendulum length will increase its period for small oscillations.
\end{hypothesisbox}
\subsection{Variables / Controls}
{\small\texttt{variablesbox}: List independent, dependent, and controlled variables.}
\begin{variablesbox}
Independent: length. Dependent: period. Controlled: release angle.
\end{variablesbox}
\subsection{Equipment / Materials}
{\small\texttt{equipmentbox}: List apparatus, quantities, and relevant specifications.}
\begin{equipmentbox}
String, a pendulum bob, a meter ruler, and a stopwatch.
\end{equipmentbox}
\subsection{Procedure / Method}
{\small\texttt{procedurebox}: Numbered steps for carrying out a repeatable task.}
\begin{procedurebox}
Measure the length, time several oscillations, and repeat the measurement.
\end{procedurebox}
\subsection{Observation}
{\small\texttt{observationbox}: Describe what was observed without mixing in interpretation.}
\begin{observationbox}
The measured period increased when the string was lengthened.
\end{observationbox}
\subsection{Data}
{\small\texttt{databox}: Measurements or a small table with units and labels.}
\begin{databox}
Record the measured length and period for every trial.
\end{databox}
\subsection{Uncertainty / Error}
{\small\texttt{uncertaintybox}: Discuss measurement uncertainty and systematic or random errors.}
\begin{uncertaintybox}
Reaction time contributes uncertainty to stopwatch measurements.
\end{uncertaintybox}
\subsection{Analysis}
{\small\texttt{analysisbox}: Explain how data or evidence supports your interpretation.}
\begin{analysisbox}
Compare the trend with the model prediction before drawing a conclusion.
\end{analysisbox}
\subsection{Conclusion}
{\small\texttt{conclusionbox}: State the result, evidence, and limits of the conclusion.}
\begin{conclusionbox}
The observed trend supports the prediction within the tested range.
\end{conclusionbox}

\clearpage
\section{Humanities, languages, and social sciences}
\subsection{Claim / Thesis}
{\small\texttt{claimbox}: A position that needs support from evidence.}
\begin{claimbox}
The author uses setting to establish the central conflict.
\end{claimbox}
\subsection{Evidence / Source}
{\small\texttt{evidencebox}: Evidence with enough source information to find it again.}
\begin{evidencebox}
Record the author, work, page or section, and the supporting detail.
\end{evidencebox}
\subsection{Quotation}
{\small\texttt{quotationbox}: An exact quotation with its source and page or location.}
\begin{quotationbox}
Place a short exact quotation here and record its source.
\end{quotationbox}
\subsection{Interpretation}
{\small\texttt{interpretationbox}: Explain what a passage, event, or piece of evidence means.}
\begin{interpretationbox}
Explain how the quoted language supports the claim.
\end{interpretationbox}
\subsection{Counterargument}
{\small\texttt{counterargumentbox}: Present an objection fairly and evaluate a response.}
\begin{counterargumentbox}
An alternative interpretation may explain the same evidence differently.
\end{counterargumentbox}
\subsection{Comparison}
{\small\texttt{comparisonbox}: Compare related ideas using the same criteria.}
\begin{comparisonbox}
Compare two theories by assumptions, predictions, evidence, and limitations.
\end{comparisonbox}
\subsection{Timeline}
{\small\texttt{timelinebox}: Events in chronological order with their significance.}
\begin{timelinebox}
Record each date, event, and why it matters.
\end{timelinebox}
\subsection{Case Study}
{\small\texttt{casebox}: Apply a model to a specific person, organization, event, or scenario.}
\begin{casebox}
Describe the situation, identify the relevant theory, and assess its fit.
\end{casebox}
\subsection{Vocabulary}
{\small\texttt{vocabularybox}: A word, meaning, pronunciation, and an example sentence.}
\begin{vocabularybox}
Word: concise. Meaning: brief but clear. Example: Write a concise summary.
\end{vocabularybox}
\subsection{Grammar / Language Pattern}
{\small\texttt{grammarbox}: A language pattern with examples and exceptions.}
\begin{grammarbox}
Record the rule, an example, and an exception if one applies.
\end{grammarbox}

\clearpage
\section{Programming, engineering, and practical work}
\subsection{Code}
{\small\texttt{codebox}: A short code listing with its language and purpose.}
\begin{codebox}
Use a lstlisting environment inside this box for literal code.
\end{codebox}
\subsection{Expected / Actual Output}
{\small\texttt{outputbox}: Record a program result, calculation output, or observed behavior.}
\begin{outputbox}
Expected output: 4.0. Actual output: 4.0.
\end{outputbox}
\subsection{Error / Fix}
{\small\texttt{debugbox}: Record the symptom, root cause, correction, and verification.}
\begin{debugbox}
A null value caused the error; validate input before dereferencing it.
\end{debugbox}
\subsection{Design Decision}
{\small\texttt{designbox}: Record a decision, alternatives, reasons, and trade-offs.}
\begin{designbox}
Separate parsing from rendering so either can change independently.
\end{designbox}
\subsection{Test Case}
{\small\texttt{testbox}: Specify input, expected result, and pass or fail evidence.}
\begin{testbox}
Input: zero net force. Expected acceleration: zero.
\end{testbox}
\subsection{Figure / Diagram}
{\small\texttt{figurebox}: A labeled diagram or image with a descriptive caption.}
\begin{figurebox}
Insert a diagram and explain what its arrows, axes, or labels represent.
\end{figurebox}
\subsection{Technique / Skill}
{\small\texttt{techniquebox}: Describe a practical technique and its success criteria.}
\begin{techniquebox}
Record the setup, movement or action, and signs of correct execution.
\end{techniquebox}
\subsection{Feedback / Critique}
{\small\texttt{feedbackbox}: Record feedback and a concrete improvement to make.}
\begin{feedbackbox}
Feedback: explain assumptions. Next revision: add an assumptions box.
\end{feedbackbox}

\clearpage
\section{Safety and cautions}
\subsection{Warning}
{\small\texttt{warning}: A caution requiring attention; yellow text on a dark background.}
\begin{warning}
Verify the operating limits before using the equipment.
\end{warning}
\subsection{Danger}
{\small\texttt{danger}: A serious hazard; red frame and title bar with a pale red body.}
\begin{danger}
Disconnect power before touching exposed electrical components.
\end{danger}

\clearpage
\section{Richer contents and customization}
\begin{formulabox}[title={Newton's Second Law},label={box:newton}]
\[
  \sum \vec{F}=m\vec{a}
\]
\NotesField{Symbols}{$m$ is mass; $\vec{a}$ is acceleration.}
\NotesField{Units}{Force in newtons, mass in kilograms, acceleration in $\mathrm{m/s^2}$.}
\NotesField{When it applies}{Constant mass, in an inertial reference frame.}
\end{formulabox}

\begin{procedurebox}
\begin{enumerate}
  \item Draw and label the forces.
  \item Choose coordinate axes.
  \item Resolve forces into components.
  \item Apply the equation along each axis.
\end{enumerate}
\end{procedurebox}

\begin{databox}[title={Measurements}]
\centering
\begin{tabular}{rr}
\toprule
Mass ($\mathrm{kg}$) & Force ($\mathrm{N}$) \\
\midrule
2 & 8 \\
5 & 20 \\
\bottomrule
\end{tabular}
\end{databox}

\begin{codebox}[title={Python: acceleration}]
\begin{lstlisting}[style=notescode,language=Python]
def acceleration(net_force, mass):
    if mass <= 0:
        raise ValueError("mass must be positive")
    return net_force / mass

print(acceleration(20, 5))
\end{lstlisting}
\end{codebox}

\begin{outputbox}
\begin{lstlisting}[style=notescode]
4.0
\end{lstlisting}
\end{outputbox}

\begin{question}[title={Question 12: custom display title}]
This changes only the displayed title, not the automatic counter.
\end{question}
\begin{answer}
The prose is italic; mathematical notation such as $F=ma$ keeps its normal mathematical formatting.
\end{answer}

\begin{note}[title={A customized title},colback=green!10!white]
Use optional keys to change one box without changing all boxes.
\end{note}

\end{document}
\fi

% See class_notes_boxes_guide.md for the complete catalog and instructions.

\ifdefined\NotesBoxesOnly
\else
  \documentclass[11pt]{article}
  \usepackage[margin=1in]{geometry}
  \usepackage[T1]{fontenc}
  \usepackage{lmodern}
\fi

\usepackage{amsmath,amssymb}
\usepackage[most]{tcolorbox}
\usepackage{graphicx}
\usepackage{booktabs}
\usepackage{listings}

% Shared appearance. The title text is always bold; body text is normal
% unless an environment (answer) explicitly changes it.
\tcbset{
  notesbase/.style={enhanced,breakable,boxrule=0.7pt,arc=2mm,
    left=3mm,right=3mm,top=2mm,bottom=2mm,
    before skip=10pt,after skip=10pt,
    fonttitle=\bfseries,fontupper=\normalfont,
    coltext=black,coltitle=white},
  notesblue/.style={colback=blue!3!white,colframe=blue!45!black,
    colbacktitle=blue!45!black},
  notesgreen/.style={colback=green!6!white,colframe=green!30!black,
    colbacktitle=green!12!white,coltitle=green!20!black},
  notesgray/.style={colback=black!2!white,colframe=black!55,
    colbacktitle=black!65},
  notesamber/.style={colback=yellow!12!white,colframe=orange!60!black,
    colbacktitle=yellow!28!white,coltitle=black!85},
  notespurple/.style={colback=violet!4!white,colframe=violet!45!black,
    colbacktitle=violet!45!black},
  noteswarning/.style={colback=black!85,colframe=yellow!80!orange,
    colbacktitle=black!90,coltitle=yellow!85!white,coltext=yellow!85!white},
  notesred/.style={colback=red!5!white,colframe=red!65!black,
    colbacktitle=red!65!black}
}

% Convenient optional fields for any box.
\newcommand{\NotesField}[2]{\par\noindent\textbf{#1:} #2\par}

% A literal code style: spaces, braces, underscores and percent signs
% inside lstlisting do not need LaTeX escaping.
\lstdefinestyle{notescode}{basicstyle=\small\ttfamily,
  breaklines=true,columns=fullflexible,keepspaces=true,
  showstringspaces=false,tabsize=4}

% BOX DEFINITIONS. All boxes accept optional tcolorbox keys in [...].

% Everyday learning
\newtcolorbox{note}[1][]{notesbase,notesgreen,
  title={Note},#1}
\newtcolorbox{definitionbox}[1][]{notesbase,notesblue,
  title={Definition},#1}
\newtcolorbox{conceptbox}[1][]{notesbase,notesblue,
  title={Key Concept},#1}
\newtcolorbox{notationbox}[1][]{notesbase,notesblue,
  title={Notation},#1}
\newtcolorbox{objectivebox}[1][]{notesbase,notesgreen,
  title={Learning Objectives},#1}
\newtcolorbox{contextbox}[1][]{notesbase,notesgray,
  title={Background / Context},#1}
\newtcolorbox{connectionbox}[1][]{notesbase,notesgreen,
  title={Connection},#1}
\newtcolorbox{summarybox}[1][]{notesbase,notesgreen,
  title={Summary},#1}

% Questions and study
\newtcolorbox[auto counter]{question}[1][]{notesbase,notesblue,
  title={Question~\thetcbcounter},#1}
\newtcolorbox{answer}[1][]{notesbase,notesgray,
  title={Answer},fontupper=\itshape,#1}
\newtcolorbox{examplebox}[1][]{notesbase,notesgray,
  title={Example},#1}
\newtcolorbox{solutionbox}[1][]{notesbase,notesgray,
  title={Worked Solution},#1}
\newtcolorbox{practicebox}[1][]{notesbase,notesblue,
  title={Practice},#1}
\newtcolorbox{recallbox}[1][]{notesbase,notespurple,
  title={Active Recall},#1}
\newtcolorbox{hintbox}[1][]{notesbase,notesgreen,
  title={Hint},#1}
\newtcolorbox{misconceptionbox}[1][]{notesbase,notesamber,
  title={Common Mistake},#1}
\newtcolorbox{askbox}[1][]{notesbase,notesamber,
  title={Ask / Clarify},#1}
\newtcolorbox{actionbox}[1][]{notesbase,notespurple,
  title={Action Item},#1}
\newtcolorbox{reflectionbox}[1][]{notesbase,notespurple,
  title={Reflection},#1}

% Math and quantitative reasoning
\newtcolorbox{formulabox}[1][]{notesbase,notesblue,
  title={Formula / Rule},#1}
\newtcolorbox{theorembox}[1][]{notesbase,notesblue,
  title={Theorem / Principle},#1}
\newtcolorbox{proofbox}[1][]{notesbase,notesgray,
  title={Proof},#1}
\newtcolorbox{derivationbox}[1][]{notesbase,notesgray,
  title={Derivation},#1}
\newtcolorbox{assumptionbox}[1][]{notesbase,notesamber,
  title={Assumptions},#1}
\newtcolorbox{givenbox}[1][]{notesbase,notesgray,
  title={Given / Find},#1}
\newtcolorbox{unitsbox}[1][]{notesbase,notesblue,
  title={Units / Dimensions},#1}
\newtcolorbox{checkbox}[1][]{notesbase,notesgreen,
  title={Result Check},#1}
\newtcolorbox{specialcasebox}[1][]{notesbase,notesamber,
  title={Special Case / Exception},#1}

% Science, labs, and research
\newtcolorbox{hypothesisbox}[1][]{notesbase,notesblue,
  title={Hypothesis},#1}
\newtcolorbox{variablesbox}[1][]{notesbase,notesblue,
  title={Variables / Controls},#1}
\newtcolorbox{equipmentbox}[1][]{notesbase,notesgray,
  title={Equipment / Materials},#1}
\newtcolorbox{procedurebox}[1][]{notesbase,notesgray,
  title={Procedure / Method},#1}
\newtcolorbox{observationbox}[1][]{notesbase,notesgray,
  title={Observation},#1}
\newtcolorbox{databox}[1][]{notesbase,notesgray,
  title={Data},#1}
\newtcolorbox{uncertaintybox}[1][]{notesbase,notesamber,
  title={Uncertainty / Error},#1}
\newtcolorbox{analysisbox}[1][]{notesbase,notesgray,
  title={Analysis},#1}
\newtcolorbox{conclusionbox}[1][]{notesbase,notesgreen,
  title={Conclusion},#1}

% Humanities, languages, and social sciences
\newtcolorbox{claimbox}[1][]{notesbase,notesblue,
  title={Claim / Thesis},#1}
\newtcolorbox{evidencebox}[1][]{notesbase,notesgray,
  title={Evidence / Source},#1}
\newtcolorbox{quotationbox}[1][]{notesbase,notesgray,
  title={Quotation},#1}
\newtcolorbox{interpretationbox}[1][]{notesbase,notespurple,
  title={Interpretation},#1}
\newtcolorbox{counterargumentbox}[1][]{notesbase,notesamber,
  title={Counterargument},#1}
\newtcolorbox{comparisonbox}[1][]{notesbase,notesgray,
  title={Comparison},#1}
\newtcolorbox{timelinebox}[1][]{notesbase,notesgray,
  title={Timeline},#1}
\newtcolorbox{casebox}[1][]{notesbase,notesgray,
  title={Case Study},#1}
\newtcolorbox{vocabularybox}[1][]{notesbase,notesblue,
  title={Vocabulary},#1}
\newtcolorbox{grammarbox}[1][]{notesbase,notesblue,
  title={Grammar / Language Pattern},#1}

% Programming, engineering, and practical work
\newtcolorbox{codebox}[1][]{notesbase,notesgray,
  title={Code},#1}
\newtcolorbox{outputbox}[1][]{notesbase,notesgray,
  title={Expected / Actual Output},#1}
\newtcolorbox{debugbox}[1][]{notesbase,notesamber,
  title={Error / Fix},#1}
\newtcolorbox{designbox}[1][]{notesbase,notespurple,
  title={Design Decision},#1}
\newtcolorbox{testbox}[1][]{notesbase,notesgreen,
  title={Test Case},#1}
\newtcolorbox{figurebox}[1][]{notesbase,notesgray,
  title={Figure / Diagram},#1}
\newtcolorbox{techniquebox}[1][]{notesbase,notesgray,
  title={Technique / Skill},#1}
\newtcolorbox{feedbackbox}[1][]{notesbase,notespurple,
  title={Feedback / Critique},#1}

% Safety and cautions
\newtcolorbox{warning}[1][]{notesbase,noteswarning,
  title={Warning},#1}
\newtcolorbox{danger}[1][]{notesbase,notesred,
  title={Danger},#1}

% DEMONSTRATION DOCUMENT: skipped when NotesBoxesOnly is defined.
\ifdefined\NotesBoxesOnly
\else
\usepackage[hidelinks]{hyperref}
\title{Class Notes Box Collection}
\author{Reusable LaTeX environments for learning and review}
\date{}
\begin{document}
\maketitle
This file contains 57 reusable box types for a broad range of classes.
The examples demonstrate formatting; replace their contents with your notes.
Use ordinary paragraphs between boxes to keep your notes readable.
\tableofcontents
\clearpage
\section{Everyday learning}
\subsection{Note}
{\small\texttt{note}: Supplementary explanation or a helpful reminder.}
\begin{note}
Review this topic after the next lecture.
\end{note}
\subsection{Definition}
{\small\texttt{definitionbox}: A term with its precise meaning.}
\begin{definitionbox}
Velocity is the rate of change of position with respect to time.
\end{definitionbox}
\subsection{Key Concept}
{\small\texttt{conceptbox}: An important idea in your own words.}
\begin{conceptbox}
A net force changes the velocity of an object.
\end{conceptbox}
\subsection{Notation}
{\small\texttt{notationbox}: Explain symbols, abbreviations, and conventions.}
\begin{notationbox}
$m$ denotes mass; $a$ denotes acceleration.
\end{notationbox}
\subsection{Learning Objectives}
{\small\texttt{objectivebox}: State what you should be able to explain or do.}
\begin{objectivebox}
Explain the difference between speed and velocity.
\end{objectivebox}
\subsection{Background / Context}
{\small\texttt{contextbox}: Give background needed to understand a topic.}
\begin{contextbox}
This lesson assumes familiarity with vectors.
\end{contextbox}
\subsection{Connection}
{\small\texttt{connectionbox}: Connect a topic to earlier lessons or another subject.}
\begin{connectionbox}
Velocity connects calculus derivatives with physical motion.
\end{connectionbox}
\subsection{Summary}
{\small\texttt{summarybox}: Record the main takeaways of a lesson.}
\begin{summarybox}
Identify forces, choose axes, and then apply the equations of motion.
\end{summarybox}

\clearpage
\section{Questions and study}
\subsection{Question}
{\small\texttt{question}: A numbered question; numbering increases automatically.}
\begin{question}
A net force of $20\,\mathrm{N}$ acts on a $5\,\mathrm{kg}$ object. Find its acceleration.
\end{question}
\subsection{Answer}
{\small\texttt{answer}: An answer with automatically italicized prose.}
\begin{answer}
$a=F_{\mathrm{net}}/m=4\,\mathrm{m/s^2}$.
\end{answer}
\subsection{Example}
{\small\texttt{examplebox}: A concrete illustration of an idea.}
\begin{examplebox}
A car travelling at constant speed around a bend still accelerates.
\end{examplebox}
\subsection{Worked Solution}
{\small\texttt{solutionbox}: Show the reasoning and steps of a calculation.}
\begin{solutionbox}
Choose the positive direction. Calculate the net force. Divide by mass.
\end{solutionbox}
\subsection{Practice}
{\small\texttt{practicebox}: An exercise to complete independently.}
\begin{practicebox}
Repeat the calculation with twice the mass.
\end{practicebox}
\subsection{Active Recall}
{\small\texttt{recallbox}: A prompt to answer without consulting the notes.}
\begin{recallbox}
State the three laws of motion in your own words.
\end{recallbox}
\subsection{Hint}
{\small\texttt{hintbox}: A small clue that does not reveal the entire solution.}
\begin{hintbox}
Start by drawing a free-body diagram.
\end{hintbox}
\subsection{Common Mistake}
{\small\texttt{misconceptionbox}: Identify an error, explain why it fails, and give a correction.}
\begin{misconceptionbox}
Constant speed does not always mean zero acceleration.
\end{misconceptionbox}
\subsection{Ask / Clarify}
{\small\texttt{askbox}: An unanswered question for an instructor or later research.}
\begin{askbox}
Under what conditions can friction be ignored?
\end{askbox}
\subsection{Action Item}
{\small\texttt{actionbox}: A task, its next step, and an optional due date.}
\begin{actionbox}
Complete practice problems 1--5 before the next class.
\end{actionbox}
\subsection{Reflection}
{\small\texttt{reflectionbox}: Evaluate your understanding, approach, or learning progress.}
\begin{reflectionbox}
I can draw the forces, but I need more practice choosing components.
\end{reflectionbox}

\clearpage
\section{Math and quantitative reasoning}
\subsection{Formula / Rule}
{\small\texttt{formulabox}: An equation or rule, with variables, units, and conditions of use.}
\begin{formulabox}
$F_{\mathrm{net}}=ma$ for constant mass in an inertial reference frame.
\end{formulabox}
\subsection{Theorem / Principle}
{\small\texttt{theorembox}: A formal result together with its hypotheses.}
\begin{theorembox}
If a function is continuous on a closed interval, it attains a maximum and a minimum.
\end{theorembox}
\subsection{Proof}
{\small\texttt{proofbox}: A logical argument establishing a claim.}
\begin{proofbox}
State the assumptions, justify each inference, and establish the conclusion.
\end{proofbox}
\subsection{Derivation}
{\small\texttt{derivationbox}: Develop an equation from definitions or established relationships.}
\begin{derivationbox}
Starting with $v=dx/dt$, integrate to relate displacement to velocity.
\end{derivationbox}
\subsection{Assumptions}
{\small\texttt{assumptionbox}: Make the conditions of a model or calculation explicit.}
\begin{assumptionbox}
Neglect air resistance and treat gravitational acceleration as constant.
\end{assumptionbox}
\subsection{Given / Find}
{\small\texttt{givenbox}: Separate known quantities from the target quantity.}
\begin{givenbox}
Given: mass and net force. Find: acceleration.
\end{givenbox}
\subsection{Units / Dimensions}
{\small\texttt{unitsbox}: Define units, convert them, or check dimensions.}
\begin{unitsbox}
$1\,\mathrm{N}=1\,\mathrm{kg\,m/s^2}$.
\end{unitsbox}
\subsection{Result Check}
{\small\texttt{checkbox}: Check sign, magnitude, units, limiting behavior, or plausibility.}
\begin{checkbox}
Doubling the mass at fixed force should halve the acceleration.
\end{checkbox}
\subsection{Special Case / Exception}
{\small\texttt{specialcasebox}: Identify boundaries, exceptions, or cases requiring different treatment.}
\begin{specialcasebox}
The expression $1/x$ is undefined at $x=0$.
\end{specialcasebox}

\clearpage
\section{Science, labs, and research}
\subsection{Hypothesis}
{\small\texttt{hypothesisbox}: A testable prediction with a rationale.}
\begin{hypothesisbox}
Increasing pendulum length will increase its period for small oscillations.
\end{hypothesisbox}
\subsection{Variables / Controls}
{\small\texttt{variablesbox}: List independent, dependent, and controlled variables.}
\begin{variablesbox}
Independent: length. Dependent: period. Controlled: release angle.
\end{variablesbox}
\subsection{Equipment / Materials}
{\small\texttt{equipmentbox}: List apparatus, quantities, and relevant specifications.}
\begin{equipmentbox}
String, a pendulum bob, a meter ruler, and a stopwatch.
\end{equipmentbox}
\subsection{Procedure / Method}
{\small\texttt{procedurebox}: Numbered steps for carrying out a repeatable task.}
\begin{procedurebox}
Measure the length, time several oscillations, and repeat the measurement.
\end{procedurebox}
\subsection{Observation}
{\small\texttt{observationbox}: Describe what was observed without mixing in interpretation.}
\begin{observationbox}
The measured period increased when the string was lengthened.
\end{observationbox}
\subsection{Data}
{\small\texttt{databox}: Measurements or a small table with units and labels.}
\begin{databox}
Record the measured length and period for every trial.
\end{databox}
\subsection{Uncertainty / Error}
{\small\texttt{uncertaintybox}: Discuss measurement uncertainty and systematic or random errors.}
\begin{uncertaintybox}
Reaction time contributes uncertainty to stopwatch measurements.
\end{uncertaintybox}
\subsection{Analysis}
{\small\texttt{analysisbox}: Explain how data or evidence supports your interpretation.}
\begin{analysisbox}
Compare the trend with the model prediction before drawing a conclusion.
\end{analysisbox}
\subsection{Conclusion}
{\small\texttt{conclusionbox}: State the result, evidence, and limits of the conclusion.}
\begin{conclusionbox}
The observed trend supports the prediction within the tested range.
\end{conclusionbox}

\clearpage
\section{Humanities, languages, and social sciences}
\subsection{Claim / Thesis}
{\small\texttt{claimbox}: A position that needs support from evidence.}
\begin{claimbox}
The author uses setting to establish the central conflict.
\end{claimbox}
\subsection{Evidence / Source}
{\small\texttt{evidencebox}: Evidence with enough source information to find it again.}
\begin{evidencebox}
Record the author, work, page or section, and the supporting detail.
\end{evidencebox}
\subsection{Quotation}
{\small\texttt{quotationbox}: An exact quotation with its source and page or location.}
\begin{quotationbox}
Place a short exact quotation here and record its source.
\end{quotationbox}
\subsection{Interpretation}
{\small\texttt{interpretationbox}: Explain what a passage, event, or piece of evidence means.}
\begin{interpretationbox}
Explain how the quoted language supports the claim.
\end{interpretationbox}
\subsection{Counterargument}
{\small\texttt{counterargumentbox}: Present an objection fairly and evaluate a response.}
\begin{counterargumentbox}
An alternative interpretation may explain the same evidence differently.
\end{counterargumentbox}
\subsection{Comparison}
{\small\texttt{comparisonbox}: Compare related ideas using the same criteria.}
\begin{comparisonbox}
Compare two theories by assumptions, predictions, evidence, and limitations.
\end{comparisonbox}
\subsection{Timeline}
{\small\texttt{timelinebox}: Events in chronological order with their significance.}
\begin{timelinebox}
Record each date, event, and why it matters.
\end{timelinebox}
\subsection{Case Study}
{\small\texttt{casebox}: Apply a model to a specific person, organization, event, or scenario.}
\begin{casebox}
Describe the situation, identify the relevant theory, and assess its fit.
\end{casebox}
\subsection{Vocabulary}
{\small\texttt{vocabularybox}: A word, meaning, pronunciation, and an example sentence.}
\begin{vocabularybox}
Word: concise. Meaning: brief but clear. Example: Write a concise summary.
\end{vocabularybox}
\subsection{Grammar / Language Pattern}
{\small\texttt{grammarbox}: A language pattern with examples and exceptions.}
\begin{grammarbox}
Record the rule, an example, and an exception if one applies.
\end{grammarbox}

\clearpage
\section{Programming, engineering, and practical work}
\subsection{Code}
{\small\texttt{codebox}: A short code listing with its language and purpose.}
\begin{codebox}
Use a lstlisting environment inside this box for literal code.
\end{codebox}
\subsection{Expected / Actual Output}
{\small\texttt{outputbox}: Record a program result, calculation output, or observed behavior.}
\begin{outputbox}
Expected output: 4.0. Actual output: 4.0.
\end{outputbox}
\subsection{Error / Fix}
{\small\texttt{debugbox}: Record the symptom, root cause, correction, and verification.}
\begin{debugbox}
A null value caused the error; validate input before dereferencing it.
\end{debugbox}
\subsection{Design Decision}
{\small\texttt{designbox}: Record a decision, alternatives, reasons, and trade-offs.}
\begin{designbox}
Separate parsing from rendering so either can change independently.
\end{designbox}
\subsection{Test Case}
{\small\texttt{testbox}: Specify input, expected result, and pass or fail evidence.}
\begin{testbox}
Input: zero net force. Expected acceleration: zero.
\end{testbox}
\subsection{Figure / Diagram}
{\small\texttt{figurebox}: A labeled diagram or image with a descriptive caption.}
\begin{figurebox}
Insert a diagram and explain what its arrows, axes, or labels represent.
\end{figurebox}
\subsection{Technique / Skill}
{\small\texttt{techniquebox}: Describe a practical technique and its success criteria.}
\begin{techniquebox}
Record the setup, movement or action, and signs of correct execution.
\end{techniquebox}
\subsection{Feedback / Critique}
{\small\texttt{feedbackbox}: Record feedback and a concrete improvement to make.}
\begin{feedbackbox}
Feedback: explain assumptions. Next revision: add an assumptions box.
\end{feedbackbox}

\clearpage
\section{Safety and cautions}
\subsection{Warning}
{\small\texttt{warning}: A caution requiring attention; yellow text on a dark background.}
\begin{warning}
Verify the operating limits before using the equipment.
\end{warning}
\subsection{Danger}
{\small\texttt{danger}: A serious hazard; red frame and title bar with a pale red body.}
\begin{danger}
Disconnect power before touching exposed electrical components.
\end{danger}

\clearpage
\section{Richer contents and customization}
\begin{formulabox}[title={Newton's Second Law},label={box:newton}]
\[
  \sum \vec{F}=m\vec{a}
\]
\NotesField{Symbols}{$m$ is mass; $\vec{a}$ is acceleration.}
\NotesField{Units}{Force in newtons, mass in kilograms, acceleration in $\mathrm{m/s^2}$.}
\NotesField{When it applies}{Constant mass, in an inertial reference frame.}
\end{formulabox}

\begin{procedurebox}
\begin{enumerate}
  \item Draw and label the forces.
  \item Choose coordinate axes.
  \item Resolve forces into components.
  \item Apply the equation along each axis.
\end{enumerate}
\end{procedurebox}

\begin{databox}[title={Measurements}]
\centering
\begin{tabular}{rr}
\toprule
Mass ($\mathrm{kg}$) & Force ($\mathrm{N}$) \\
\midrule
2 & 8 \\
5 & 20 \\
\bottomrule
\end{tabular}
\end{databox}

\begin{codebox}[title={Python: acceleration}]
\begin{lstlisting}[style=notescode,language=Python]
def acceleration(net_force, mass):
    if mass <= 0:
        raise ValueError("mass must be positive")
    return net_force / mass

print(acceleration(20, 5))
\end{lstlisting}
\end{codebox}

\begin{outputbox}
\begin{lstlisting}[style=notescode]
4.0
\end{lstlisting}
\end{outputbox}

\begin{question}[title={Question 12: custom display title}]
This changes only the displayed title, not the automatic counter.
\end{question}
\begin{answer}
The prose is italic; mathematical notation such as $F=ma$ keeps its normal mathematical formatting.
\end{answer}

\begin{note}[title={A customized title},colback=green!10!white]
Use optional keys to change one box without changing all boxes.
\end{note}

\end{document}
\fi

% See class_notes_boxes_guide.md for the complete catalog and instructions.

\ifdefined\NotesBoxesOnly
\else
  \documentclass[11pt]{article}
  \usepackage[margin=1in]{geometry}
  \usepackage[T1]{fontenc}
  \usepackage{lmodern}
\fi

\usepackage{amsmath,amssymb}
\usepackage[most]{tcolorbox}
\usepackage{graphicx}
\usepackage{booktabs}
\usepackage{listings}

% Shared appearance. The title text is always bold; body text is normal
% unless an environment (answer) explicitly changes it.
\tcbset{
  notesbase/.style={enhanced,breakable,boxrule=0.7pt,arc=2mm,
    left=3mm,right=3mm,top=2mm,bottom=2mm,
    before skip=10pt,after skip=10pt,
    fonttitle=\bfseries,fontupper=\normalfont,
    coltext=black,coltitle=white},
  notesblue/.style={colback=blue!3!white,colframe=blue!45!black,
    colbacktitle=blue!45!black},
  notesgreen/.style={colback=green!6!white,colframe=green!30!black,
    colbacktitle=green!12!white,coltitle=green!20!black},
  notesgray/.style={colback=black!2!white,colframe=black!55,
    colbacktitle=black!65},
  notesamber/.style={colback=yellow!12!white,colframe=orange!60!black,
    colbacktitle=yellow!28!white,coltitle=black!85},
  notespurple/.style={colback=violet!4!white,colframe=violet!45!black,
    colbacktitle=violet!45!black},
  noteswarning/.style={colback=black!85,colframe=yellow!80!orange,
    colbacktitle=black!90,coltitle=yellow!85!white,coltext=yellow!85!white},
  notesred/.style={colback=red!5!white,colframe=red!65!black,
    colbacktitle=red!65!black}
}

% Convenient optional fields for any box.
\newcommand{\NotesField}[2]{\par\noindent\textbf{#1:} #2\par}

% A literal code style: spaces, braces, underscores and percent signs
% inside lstlisting do not need LaTeX escaping.
\lstdefinestyle{notescode}{basicstyle=\small\ttfamily,
  breaklines=true,columns=fullflexible,keepspaces=true,
  showstringspaces=false,tabsize=4}

% BOX DEFINITIONS. All boxes accept optional tcolorbox keys in [...].

% Everyday learning
\newtcolorbox{note}[1][]{notesbase,notesgreen,
  title={Note},#1}
\newtcolorbox{definitionbox}[1][]{notesbase,notesblue,
  title={Definition},#1}
\newtcolorbox{conceptbox}[1][]{notesbase,notesblue,
  title={Key Concept},#1}
\newtcolorbox{notationbox}[1][]{notesbase,notesblue,
  title={Notation},#1}
\newtcolorbox{objectivebox}[1][]{notesbase,notesgreen,
  title={Learning Objectives},#1}
\newtcolorbox{contextbox}[1][]{notesbase,notesgray,
  title={Background / Context},#1}
\newtcolorbox{connectionbox}[1][]{notesbase,notesgreen,
  title={Connection},#1}
\newtcolorbox{summarybox}[1][]{notesbase,notesgreen,
  title={Summary},#1}

% Questions and study
\newtcolorbox[auto counter]{question}[1][]{notesbase,notesblue,
  title={Question~\thetcbcounter},#1}
\newtcolorbox{answer}[1][]{notesbase,notesgray,
  title={Answer},fontupper=\itshape,#1}
\newtcolorbox{examplebox}[1][]{notesbase,notesgray,
  title={Example},#1}
\newtcolorbox{solutionbox}[1][]{notesbase,notesgray,
  title={Worked Solution},#1}
\newtcolorbox{practicebox}[1][]{notesbase,notesblue,
  title={Practice},#1}
\newtcolorbox{recallbox}[1][]{notesbase,notespurple,
  title={Active Recall},#1}
\newtcolorbox{hintbox}[1][]{notesbase,notesgreen,
  title={Hint},#1}
\newtcolorbox{misconceptionbox}[1][]{notesbase,notesamber,
  title={Common Mistake},#1}
\newtcolorbox{askbox}[1][]{notesbase,notesamber,
  title={Ask / Clarify},#1}
\newtcolorbox{actionbox}[1][]{notesbase,notespurple,
  title={Action Item},#1}
\newtcolorbox{reflectionbox}[1][]{notesbase,notespurple,
  title={Reflection},#1}

% Math and quantitative reasoning
\newtcolorbox{formulabox}[1][]{notesbase,notesblue,
  title={Formula / Rule},#1}
\newtcolorbox{theorembox}[1][]{notesbase,notesblue,
  title={Theorem / Principle},#1}
\newtcolorbox{proofbox}[1][]{notesbase,notesgray,
  title={Proof},#1}
\newtcolorbox{derivationbox}[1][]{notesbase,notesgray,
  title={Derivation},#1}
\newtcolorbox{assumptionbox}[1][]{notesbase,notesamber,
  title={Assumptions},#1}
\newtcolorbox{givenbox}[1][]{notesbase,notesgray,
  title={Given / Find},#1}
\newtcolorbox{unitsbox}[1][]{notesbase,notesblue,
  title={Units / Dimensions},#1}
\newtcolorbox{checkbox}[1][]{notesbase,notesgreen,
  title={Result Check},#1}
\newtcolorbox{specialcasebox}[1][]{notesbase,notesamber,
  title={Special Case / Exception},#1}

% Science, labs, and research
\newtcolorbox{hypothesisbox}[1][]{notesbase,notesblue,
  title={Hypothesis},#1}
\newtcolorbox{variablesbox}[1][]{notesbase,notesblue,
  title={Variables / Controls},#1}
\newtcolorbox{equipmentbox}[1][]{notesbase,notesgray,
  title={Equipment / Materials},#1}
\newtcolorbox{procedurebox}[1][]{notesbase,notesgray,
  title={Procedure / Method},#1}
\newtcolorbox{observationbox}[1][]{notesbase,notesgray,
  title={Observation},#1}
\newtcolorbox{databox}[1][]{notesbase,notesgray,
  title={Data},#1}
\newtcolorbox{uncertaintybox}[1][]{notesbase,notesamber,
  title={Uncertainty / Error},#1}
\newtcolorbox{analysisbox}[1][]{notesbase,notesgray,
  title={Analysis},#1}
\newtcolorbox{conclusionbox}[1][]{notesbase,notesgreen,
  title={Conclusion},#1}

% Humanities, languages, and social sciences
\newtcolorbox{claimbox}[1][]{notesbase,notesblue,
  title={Claim / Thesis},#1}
\newtcolorbox{evidencebox}[1][]{notesbase,notesgray,
  title={Evidence / Source},#1}
\newtcolorbox{quotationbox}[1][]{notesbase,notesgray,
  title={Quotation},#1}
\newtcolorbox{interpretationbox}[1][]{notesbase,notespurple,
  title={Interpretation},#1}
\newtcolorbox{counterargumentbox}[1][]{notesbase,notesamber,
  title={Counterargument},#1}
\newtcolorbox{comparisonbox}[1][]{notesbase,notesgray,
  title={Comparison},#1}
\newtcolorbox{timelinebox}[1][]{notesbase,notesgray,
  title={Timeline},#1}
\newtcolorbox{casebox}[1][]{notesbase,notesgray,
  title={Case Study},#1}
\newtcolorbox{vocabularybox}[1][]{notesbase,notesblue,
  title={Vocabulary},#1}
\newtcolorbox{grammarbox}[1][]{notesbase,notesblue,
  title={Grammar / Language Pattern},#1}

% Programming, engineering, and practical work
\newtcolorbox{codebox}[1][]{notesbase,notesgray,
  title={Code},#1}
\newtcolorbox{outputbox}[1][]{notesbase,notesgray,
  title={Expected / Actual Output},#1}
\newtcolorbox{debugbox}[1][]{notesbase,notesamber,
  title={Error / Fix},#1}
\newtcolorbox{designbox}[1][]{notesbase,notespurple,
  title={Design Decision},#1}
\newtcolorbox{testbox}[1][]{notesbase,notesgreen,
  title={Test Case},#1}
\newtcolorbox{figurebox}[1][]{notesbase,notesgray,
  title={Figure / Diagram},#1}
\newtcolorbox{techniquebox}[1][]{notesbase,notesgray,
  title={Technique / Skill},#1}
\newtcolorbox{feedbackbox}[1][]{notesbase,notespurple,
  title={Feedback / Critique},#1}

% Safety and cautions
\newtcolorbox{warning}[1][]{notesbase,noteswarning,
  title={Warning},#1}
\newtcolorbox{danger}[1][]{notesbase,notesred,
  title={Danger},#1}

% DEMONSTRATION DOCUMENT: skipped when NotesBoxesOnly is defined.
\ifdefined\NotesBoxesOnly
\else
\usepackage[hidelinks]{hyperref}
\title{Class Notes Box Collection}
\author{Reusable LaTeX environments for learning and review}
\date{}
\begin{document}
\maketitle
This file contains 57 reusable box types for a broad range of classes.
The examples demonstrate formatting; replace their contents with your notes.
Use ordinary paragraphs between boxes to keep your notes readable.
\tableofcontents
\clearpage
\section{Everyday learning}
\subsection{Note}
{\small\texttt{note}: Supplementary explanation or a helpful reminder.}
\begin{note}
Review this topic after the next lecture.
\end{note}
\subsection{Definition}
{\small\texttt{definitionbox}: A term with its precise meaning.}
\begin{definitionbox}
Velocity is the rate of change of position with respect to time.
\end{definitionbox}
\subsection{Key Concept}
{\small\texttt{conceptbox}: An important idea in your own words.}
\begin{conceptbox}
A net force changes the velocity of an object.
\end{conceptbox}
\subsection{Notation}
{\small\texttt{notationbox}: Explain symbols, abbreviations, and conventions.}
\begin{notationbox}
$m$ denotes mass; $a$ denotes acceleration.
\end{notationbox}
\subsection{Learning Objectives}
{\small\texttt{objectivebox}: State what you should be able to explain or do.}
\begin{objectivebox}
Explain the difference between speed and velocity.
\end{objectivebox}
\subsection{Background / Context}
{\small\texttt{contextbox}: Give background needed to understand a topic.}
\begin{contextbox}
This lesson assumes familiarity with vectors.
\end{contextbox}
\subsection{Connection}
{\small\texttt{connectionbox}: Connect a topic to earlier lessons or another subject.}
\begin{connectionbox}
Velocity connects calculus derivatives with physical motion.
\end{connectionbox}
\subsection{Summary}
{\small\texttt{summarybox}: Record the main takeaways of a lesson.}
\begin{summarybox}
Identify forces, choose axes, and then apply the equations of motion.
\end{summarybox}

\clearpage
\section{Questions and study}
\subsection{Question}
{\small\texttt{question}: A numbered question; numbering increases automatically.}
\begin{question}
A net force of $20\,\mathrm{N}$ acts on a $5\,\mathrm{kg}$ object. Find its acceleration.
\end{question}
\subsection{Answer}
{\small\texttt{answer}: An answer with automatically italicized prose.}
\begin{answer}
$a=F_{\mathrm{net}}/m=4\,\mathrm{m/s^2}$.
\end{answer}
\subsection{Example}
{\small\texttt{examplebox}: A concrete illustration of an idea.}
\begin{examplebox}
A car travelling at constant speed around a bend still accelerates.
\end{examplebox}
\subsection{Worked Solution}
{\small\texttt{solutionbox}: Show the reasoning and steps of a calculation.}
\begin{solutionbox}
Choose the positive direction. Calculate the net force. Divide by mass.
\end{solutionbox}
\subsection{Practice}
{\small\texttt{practicebox}: An exercise to complete independently.}
\begin{practicebox}
Repeat the calculation with twice the mass.
\end{practicebox}
\subsection{Active Recall}
{\small\texttt{recallbox}: A prompt to answer without consulting the notes.}
\begin{recallbox}
State the three laws of motion in your own words.
\end{recallbox}
\subsection{Hint}
{\small\texttt{hintbox}: A small clue that does not reveal the entire solution.}
\begin{hintbox}
Start by drawing a free-body diagram.
\end{hintbox}
\subsection{Common Mistake}
{\small\texttt{misconceptionbox}: Identify an error, explain why it fails, and give a correction.}
\begin{misconceptionbox}
Constant speed does not always mean zero acceleration.
\end{misconceptionbox}
\subsection{Ask / Clarify}
{\small\texttt{askbox}: An unanswered question for an instructor or later research.}
\begin{askbox}
Under what conditions can friction be ignored?
\end{askbox}
\subsection{Action Item}
{\small\texttt{actionbox}: A task, its next step, and an optional due date.}
\begin{actionbox}
Complete practice problems 1--5 before the next class.
\end{actionbox}
\subsection{Reflection}
{\small\texttt{reflectionbox}: Evaluate your understanding, approach, or learning progress.}
\begin{reflectionbox}
I can draw the forces, but I need more practice choosing components.
\end{reflectionbox}

\clearpage
\section{Math and quantitative reasoning}
\subsection{Formula / Rule}
{\small\texttt{formulabox}: An equation or rule, with variables, units, and conditions of use.}
\begin{formulabox}
$F_{\mathrm{net}}=ma$ for constant mass in an inertial reference frame.
\end{formulabox}
\subsection{Theorem / Principle}
{\small\texttt{theorembox}: A formal result together with its hypotheses.}
\begin{theorembox}
If a function is continuous on a closed interval, it attains a maximum and a minimum.
\end{theorembox}
\subsection{Proof}
{\small\texttt{proofbox}: A logical argument establishing a claim.}
\begin{proofbox}
State the assumptions, justify each inference, and establish the conclusion.
\end{proofbox}
\subsection{Derivation}
{\small\texttt{derivationbox}: Develop an equation from definitions or established relationships.}
\begin{derivationbox}
Starting with $v=dx/dt$, integrate to relate displacement to velocity.
\end{derivationbox}
\subsection{Assumptions}
{\small\texttt{assumptionbox}: Make the conditions of a model or calculation explicit.}
\begin{assumptionbox}
Neglect air resistance and treat gravitational acceleration as constant.
\end{assumptionbox}
\subsection{Given / Find}
{\small\texttt{givenbox}: Separate known quantities from the target quantity.}
\begin{givenbox}
Given: mass and net force. Find: acceleration.
\end{givenbox}
\subsection{Units / Dimensions}
{\small\texttt{unitsbox}: Define units, convert them, or check dimensions.}
\begin{unitsbox}
$1\,\mathrm{N}=1\,\mathrm{kg\,m/s^2}$.
\end{unitsbox}
\subsection{Result Check}
{\small\texttt{checkbox}: Check sign, magnitude, units, limiting behavior, or plausibility.}
\begin{checkbox}
Doubling the mass at fixed force should halve the acceleration.
\end{checkbox}
\subsection{Special Case / Exception}
{\small\texttt{specialcasebox}: Identify boundaries, exceptions, or cases requiring different treatment.}
\begin{specialcasebox}
The expression $1/x$ is undefined at $x=0$.
\end{specialcasebox}

\clearpage
\section{Science, labs, and research}
\subsection{Hypothesis}
{\small\texttt{hypothesisbox}: A testable prediction with a rationale.}
\begin{hypothesisbox}
Increasing pendulum length will increase its period for small oscillations.
\end{hypothesisbox}
\subsection{Variables / Controls}
{\small\texttt{variablesbox}: List independent, dependent, and controlled variables.}
\begin{variablesbox}
Independent: length. Dependent: period. Controlled: release angle.
\end{variablesbox}
\subsection{Equipment / Materials}
{\small\texttt{equipmentbox}: List apparatus, quantities, and relevant specifications.}
\begin{equipmentbox}
String, a pendulum bob, a meter ruler, and a stopwatch.
\end{equipmentbox}
\subsection{Procedure / Method}
{\small\texttt{procedurebox}: Numbered steps for carrying out a repeatable task.}
\begin{procedurebox}
Measure the length, time several oscillations, and repeat the measurement.
\end{procedurebox}
\subsection{Observation}
{\small\texttt{observationbox}: Describe what was observed without mixing in interpretation.}
\begin{observationbox}
The measured period increased when the string was lengthened.
\end{observationbox}
\subsection{Data}
{\small\texttt{databox}: Measurements or a small table with units and labels.}
\begin{databox}
Record the measured length and period for every trial.
\end{databox}
\subsection{Uncertainty / Error}
{\small\texttt{uncertaintybox}: Discuss measurement uncertainty and systematic or random errors.}
\begin{uncertaintybox}
Reaction time contributes uncertainty to stopwatch measurements.
\end{uncertaintybox}
\subsection{Analysis}
{\small\texttt{analysisbox}: Explain how data or evidence supports your interpretation.}
\begin{analysisbox}
Compare the trend with the model prediction before drawing a conclusion.
\end{analysisbox}
\subsection{Conclusion}
{\small\texttt{conclusionbox}: State the result, evidence, and limits of the conclusion.}
\begin{conclusionbox}
The observed trend supports the prediction within the tested range.
\end{conclusionbox}

\clearpage
\section{Humanities, languages, and social sciences}
\subsection{Claim / Thesis}
{\small\texttt{claimbox}: A position that needs support from evidence.}
\begin{claimbox}
The author uses setting to establish the central conflict.
\end{claimbox}
\subsection{Evidence / Source}
{\small\texttt{evidencebox}: Evidence with enough source information to find it again.}
\begin{evidencebox}
Record the author, work, page or section, and the supporting detail.
\end{evidencebox}
\subsection{Quotation}
{\small\texttt{quotationbox}: An exact quotation with its source and page or location.}
\begin{quotationbox}
Place a short exact quotation here and record its source.
\end{quotationbox}
\subsection{Interpretation}
{\small\texttt{interpretationbox}: Explain what a passage, event, or piece of evidence means.}
\begin{interpretationbox}
Explain how the quoted language supports the claim.
\end{interpretationbox}
\subsection{Counterargument}
{\small\texttt{counterargumentbox}: Present an objection fairly and evaluate a response.}
\begin{counterargumentbox}
An alternative interpretation may explain the same evidence differently.
\end{counterargumentbox}
\subsection{Comparison}
{\small\texttt{comparisonbox}: Compare related ideas using the same criteria.}
\begin{comparisonbox}
Compare two theories by assumptions, predictions, evidence, and limitations.
\end{comparisonbox}
\subsection{Timeline}
{\small\texttt{timelinebox}: Events in chronological order with their significance.}
\begin{timelinebox}
Record each date, event, and why it matters.
\end{timelinebox}
\subsection{Case Study}
{\small\texttt{casebox}: Apply a model to a specific person, organization, event, or scenario.}
\begin{casebox}
Describe the situation, identify the relevant theory, and assess its fit.
\end{casebox}
\subsection{Vocabulary}
{\small\texttt{vocabularybox}: A word, meaning, pronunciation, and an example sentence.}
\begin{vocabularybox}
Word: concise. Meaning: brief but clear. Example: Write a concise summary.
\end{vocabularybox}
\subsection{Grammar / Language Pattern}
{\small\texttt{grammarbox}: A language pattern with examples and exceptions.}
\begin{grammarbox}
Record the rule, an example, and an exception if one applies.
\end{grammarbox}

\clearpage
\section{Programming, engineering, and practical work}
\subsection{Code}
{\small\texttt{codebox}: A short code listing with its language and purpose.}
\begin{codebox}
Use a lstlisting environment inside this box for literal code.
\end{codebox}
\subsection{Expected / Actual Output}
{\small\texttt{outputbox}: Record a program result, calculation output, or observed behavior.}
\begin{outputbox}
Expected output: 4.0. Actual output: 4.0.
\end{outputbox}
\subsection{Error / Fix}
{\small\texttt{debugbox}: Record the symptom, root cause, correction, and verification.}
\begin{debugbox}
A null value caused the error; validate input before dereferencing it.
\end{debugbox}
\subsection{Design Decision}
{\small\texttt{designbox}: Record a decision, alternatives, reasons, and trade-offs.}
\begin{designbox}
Separate parsing from rendering so either can change independently.
\end{designbox}
\subsection{Test Case}
{\small\texttt{testbox}: Specify input, expected result, and pass or fail evidence.}
\begin{testbox}
Input: zero net force. Expected acceleration: zero.
\end{testbox}
\subsection{Figure / Diagram}
{\small\texttt{figurebox}: A labeled diagram or image with a descriptive caption.}
\begin{figurebox}
Insert a diagram and explain what its arrows, axes, or labels represent.
\end{figurebox}
\subsection{Technique / Skill}
{\small\texttt{techniquebox}: Describe a practical technique and its success criteria.}
\begin{techniquebox}
Record the setup, movement or action, and signs of correct execution.
\end{techniquebox}
\subsection{Feedback / Critique}
{\small\texttt{feedbackbox}: Record feedback and a concrete improvement to make.}
\begin{feedbackbox}
Feedback: explain assumptions. Next revision: add an assumptions box.
\end{feedbackbox}

\clearpage
\section{Safety and cautions}
\subsection{Warning}
{\small\texttt{warning}: A caution requiring attention; yellow text on a dark background.}
\begin{warning}
Verify the operating limits before using the equipment.
\end{warning}
\subsection{Danger}
{\small\texttt{danger}: A serious hazard; red frame and title bar with a pale red body.}
\begin{danger}
Disconnect power before touching exposed electrical components.
\end{danger}

\clearpage
\section{Richer contents and customization}
\begin{formulabox}[title={Newton's Second Law},label={box:newton}]
\[
  \sum \vec{F}=m\vec{a}
\]
\NotesField{Symbols}{$m$ is mass; $\vec{a}$ is acceleration.}
\NotesField{Units}{Force in newtons, mass in kilograms, acceleration in $\mathrm{m/s^2}$.}
\NotesField{When it applies}{Constant mass, in an inertial reference frame.}
\end{formulabox}

\begin{procedurebox}
\begin{enumerate}
  \item Draw and label the forces.
  \item Choose coordinate axes.
  \item Resolve forces into components.
  \item Apply the equation along each axis.
\end{enumerate}
\end{procedurebox}

\begin{databox}[title={Measurements}]
\centering
\begin{tabular}{rr}
\toprule
Mass ($\mathrm{kg}$) & Force ($\mathrm{N}$) \\
\midrule
2 & 8 \\
5 & 20 \\
\bottomrule
\end{tabular}
\end{databox}

\begin{codebox}[title={Python: acceleration}]
\begin{lstlisting}[style=notescode,language=Python]
def acceleration(net_force, mass):
    if mass <= 0:
        raise ValueError("mass must be positive")
    return net_force / mass

print(acceleration(20, 5))
\end{lstlisting}
\end{codebox}

\begin{outputbox}
\begin{lstlisting}[style=notescode]
4.0
\end{lstlisting}
\end{outputbox}

\begin{question}[title={Question 12: custom display title}]
This changes only the displayed title, not the automatic counter.
\end{question}
\begin{answer}
The prose is italic; mathematical notation such as $F=ma$ keeps its normal mathematical formatting.
\end{answer}

\begin{note}[title={A customized title},colback=green!10!white]
Use optional keys to change one box without changing all boxes.
\end{note}

\end{document}
\fi
